% Questions de cours

\begin{taggedblock}{1}
\begin{exoq}
Donner la définition d'une donnée.
\end{exoq} 
\begin{exor}
Une donnée est une valeur décrivant une entité (objet, personne, événement, ...). \bqr
\end{exor} 
\end{taggedblock}

\begin{taggedblock}{2}
\begin{exoq}
Qu'est-ce qu'une donnée élémentaire ?
\end{exoq} 
\begin{exor}
Une donnée élémentaire est une donnée constituée d'une seule valeur. \bqr
\end{exor} 
\end{taggedblock}

\begin{taggedblock}{3}
\begin{exoq}
Qu'est ce qu'une donnée complexe ?
\end{exoq} 
\begin{exor}
Une donnée complexe est une donnée constituée de plusieurs données élémentaires. \bqr
\end{exor} 
\end{taggedblock}

\begin{taggedblock}{4}
\begin{exoq}
Citer un descripteur dans une adresse postale.
\end{exoq} 
\begin{exor}
Le n° de voirie, le code postale, le nom de la voirie,... \bqr
\end{exor} 
\end{taggedblock}

\begin{taggedblock}{5}
\begin{exoq}
Donner les noms des deux catégories de supports de stockage.
\end{exoq} 
\begin{exor}
Les fichiers et les bases de données sont deux catégories de supports de stockage. \bqr
\end{exor} 
\end{taggedblock}

\begin{taggedblock}{6}
\begin{exoq}
Citer trois opérations du cycle de vie d'une donnée (ne pas détailler).
\end{exoq} 
\begin{exor}
Exemples d'opérations : Traitement et partage, collecte, archivage, destruction, analyse, sauvegarde. \bqr
\end{exor} 
\end{taggedblock}

\begin{taggedblock}{7}
\begin{exoq}
Donner deux exemples de typages de données.
\end{exoq} 
\begin{exor}
Exemples de typages : chaîne de caractères, nombre entier, nombre décimal, date, … \bqr
\end{exor} 
\end{taggedblock}

\begin{taggedblock}{8}
\begin{exoq}
Donner le nom de cette opération de la vie d'une donnée qui consiste à recueillir les données.
\end{exoq} 
\begin{exor}
Cette opération se nomme la collecte. \bqr
\end{exor} 
\end{taggedblock}

\begin{taggedblock}{9}
\begin{exoq}
Qu'est-ce qu'une collection ?
\end{exoq} 
\begin{exor}
Une collection regroupe des objets possédant les mêmes descripteurs. \bqr
\end{exor} 
\end{taggedblock}

\begin{taggedblock}{10}
\begin{exoq}
Citer deux types de formats de fichiers.
\end{exoq} 
\begin{exor}
Le CSV et le JSON sont deux types de formats de fichiers. \bqr
\end{exor} 
\end{taggedblock}

\begin{taggedblock}{11}
\begin{exoq}
Dans un fichier CSV, comment sont séparées les données ?
\end{exoq} 
\begin{exor}
Dans un fichier CSV, les données sont séparées par un caractère spécifique. \bqr
\end{exor} 
\end{taggedblock}

\begin{taggedblock}{12}
\begin{exoq}
Quel est le nom de la dernière opération de la vie d'une donnée ?
\end{exoq} 
\begin{exor}
La dernière opération est la destruction de la donnée. \bqr
\end{exor} 
\end{taggedblock}

\begin{taggedblock}{13}
\begin{exoq}
Donner un exemple de séparateur dans un fichier CSV.
\end{exoq} 
\begin{exor}
Dans un fichier CSV, le séparateur peut être une virgule ou un point-virgule. \bqr
\end{exor} 
\end{taggedblock}

\begin{taggedblock}{14}
\begin{exoq}
Quel est l'avantage du format JSON sur le format CSV ?
\end{exoq} 
\begin{exor}
Le format JSON permet de stocker des données plus complexes. \bqr
\end{exor} 
\end{taggedblock}

\begin{taggedblock}{15}
\begin{exoq}
Que signifie et représente k dans ko ?
\end{exoq} 
\begin{exor}
k signifie kilo et représente $ \times 10^{3} $ \bqr
\end{exor} 
\end{taggedblock}

\begin{taggedblock}{16}
\begin{exoq}
Que signifie et représente M dans Mo ?
\end{exoq} 
\begin{exor}
M signifie méga et représente $ \times 10^{6} $ \bqr
\end{exor} 
\end{taggedblock}

\begin{taggedblock}{17}
\begin{exoq}
Que signifie et représente G dans Go ?
\end{exoq} 
\begin{exor}
G signifie giga et représente $ \times 10^{9} $ \bqr
\end{exor} 
\end{taggedblock}

\begin{taggedblock}{18}
\begin{exoq}
Que signifie et représente T dans To ?
\end{exoq} 
\begin{exor}
T signifie téra et représente $ \times 10^{12} $ \bqr
\end{exor} 
\end{taggedblock}

\begin{taggedblock}{19}
\begin{exoq}
Attribuer les années (1928 1956 1982 1995 1999 2016) aux supports de stockage suivants :

clé usb, disque dur, disque SSD, CDrom, DVDrom, carte perforée.
\end{exoq} 
\begin{exor}
clé usb (1999), disque dur (1956), disque SSD (2016), CDrom (1982), DVDrom (1995), carte perforée (1928). \bqr
\end{exor} 
\end{taggedblock}

\begin{taggedblock}{20}
\begin{exoq}
En quoi consiste un traitement sur des données ?
\end{exoq} 
\begin{exor}
Un traitement sur des données consiste à effectuer des actions afin d'obtenir un résultat. \bqr
\end{exor} 
\end{taggedblock}

\begin{taggedblock}{21}
\begin{exoq}
Donner 2 exemples de traitements sur des données.
\end{exoq} 
\begin{exor}
Trier, filtrer, faire des calculs, mettre en forme sont des exemples de traitements. \bqr
\end{exor} 
\end{taggedblock}

\begin{taggedblock}{22}
\begin{exoq}
Qu'est-ce qu'une métadonnée ?
\end{exoq} 
\begin{exor}
Une métadonnée est une donnée qui sert à définir ou à décrire d'autres données principales afin de les valoriser pour mieux les rechercher. \bqr
\end{exor} 
\end{taggedblock}

\begin{taggedblock}{23}
\begin{exoq}
Donner deux exemples de métadonnée.
\end{exoq} 
\begin{exor}
Exemples de métadonnée : une date de sauvegarde ou de publication, une taille de fichier, un auteur, ... \bqr
\end{exor} 
\end{taggedblock}

\begin{taggedblock}{24}
\begin{exoq}
Qu'est-ce qu'une donnée personnelle ?
\end{exoq} 
\begin{exor}
Une donnée personnelle est une donnée qui identifie directement ou indirectement une personne physique. \bqr
\end{exor} 
\end{taggedblock}

\begin{taggedblock}{25}
\begin{exoq}
Que signifie RGPD ?
\end{exoq} 
\begin{exor}
RGPD signifie : Règlement Général de la Protection des Données. \bqr
\end{exor} 
\end{taggedblock}

\begin{taggedblock}{26}
\begin{exoq}
En France, quel organisme est chargé de veiller au respect des données privées ?
\end{exoq} 
\begin{exor}
C'est la CNIL qui se charge de veiller au respect des données privées. \bqr
\end{exor} 
\end{taggedblock}

\begin{taggedblock}{27}
\begin{exoq}
En quoi le langage Python est-il intéressant au niveau de la manipulation de données ?
\end{exoq} 
\begin{exor}
Python permet d'automatiser les traitements de données. \bqr
\end{exor} 
\end{taggedblock}

\begin{taggedblock}{28}
\begin{exoq}
Que signifie CSV ?
\end{exoq} 
\begin{exor}
CSV signifie Coma Separated Values. \bqr
\end{exor} 
\end{taggedblock}

\begin{taggedblock}{29}
\begin{exoq}
Quel type de données peuvent s'enregistrer dans un fichier CSV ?
\end{exoq} 
\begin{exor}
Un fichier CSV contient des données tabulaires (données structurées). \bqr
\end{exor} 
\end{taggedblock}

\begin{taggedblock}{30}
\begin{exoq}
Comment nomme-t-on la première ligne d'un fichier CSV ?
\end{exoq} 
\begin{exor}
La première ligne d'un fichier CSV se nomme l'entête. \bqr
\end{exor} 
\end{taggedblock}

\begin{taggedblock}{31}
\begin{exoq}
Que contient l'entête d'un fichier CSV ?
\end{exoq} 
\begin{exor}
L'entête d'un fichier CSV contient des descripteurs. \bqr
\end{exor} 
\end{taggedblock}

\begin{taggedblock}{32}
\begin{exoq}
Citer deux moyens de créer un fichier CSV.
\end{exoq} 
\begin{exor}
On peut créer un fichier CSV en utilisant un tableur comme Excel, ou en utilisant un langage de programmation comme Python. \bqr
\end{exor} 
\end{taggedblock}
